\originalchapter{怎样读 R1 的奖励、推理轨迹与证据}\label{study:r1}
\noindent 阅读对象是 2025-09-17 的 Nature 正式论文\cite{r1nature}，其完整正文、Methods 与 Extended Data 译文另卷提供。该版 CC BY 4.0 不自动适用于另一个 arXiv 版本或独立补充文件。本说明以原创推导把论文放到基础学习、DPO、测试时计算与可验证训练的共同问题链中。

\originalsection{从目标概率到奖励梯度}
预训练给出了 $\pi_\theta(y\mid x)$，自回归分解为各 token 条件概率的乘积。监督微调用固定示范的负对数概率训练；奖励训练则先从当前或旧策略采样，再依据回答是否满足目标调整分布。若输出具有固定支持，在其上 $\pi_\theta(y\mid x)>0$ 且关于 $\theta$ 可微，$r(x,y)$ 不显式依赖 $\theta$，并且可以交换求导与求和，则
\[
\nabla_\theta\E_{y\sim\pi_\theta}[r(x,y)]
=\sum_y r(x,y)\pi_\theta(y\mid x)\nabla_\theta\log\pi_\theta(y\mid x)
=\E[r(x,y)\nabla_\theta\log\pi_\theta(y\mid x)].
\]
若奖励依赖生成出的文本，通过评分器产生一个数，这个恒等式并不要求对评分程序求导。数学推导说明怎样提高有奖励输出的概率；它没有说明奖励为什么忠实地衡量了推理质量。

对任意只依赖输入 $x$ 的基线 $b(x)$，$\E[b(x)\nabla\log\pi]=b(x)\nabla\sum_y\pi(y\mid x)=0$。减基线可改变方差而不改变该精确期望。采样组均值包含当前样本，是否无偏还须看具体估计器；这正是教材分开推导组内与留一基线的原因。

\originalsection{GRPO 的两个粒度不能混算}
Nature Methods 式(1)以整条输出策略比 $\pi_\theta(o_i\mid q)/\pi_{\theta_{\rm old}}(o_i\mid q)$ 写示意目标。token 实现又可在每个时间步用比值 $\rho_{it}$ 并对裁剪项平均。整序列比是 token 比的乘积；对乘积裁剪和对每个比值分别裁剪后平均一般不相同。读者应在论文版本、实现与本地简化之间明确粒度，而不是把相似符号当成同一算法。

\begin{example}
两 token 的旧条件概率分别为 $0.4,0.5$，新条件概率为 $0.8,0.25$，则两 token 比为 $2,0.5$，整序列比为 $1$。取正优势 $A=1$、裁剪区间 $[0.8,1.2]$，序列裁剪项为 $1$；token 裁剪目标的均值为 $(1.2+0.5)/2=0.85$。
\end{example}
\begin{solution}
PPO 型项是 $\min(\rho A,\operatorname{clip}(\rho)A)$，第二个 token 的原比 $0.5$ 比裁剪后 $0.8$ 更小，故仍取 $0.5$。这不是两个公式算错，而是两种目标不同。
\end{solution}

\originalsection{可验证任务为什么提供了新的训练条件}
数学答案和程序测试让部分任务有廉价、重复可用的反馈；同一输入可生成许多轨迹，探索不同策略。这里可验证的对象可能只是最终答案或有限测试，并非整段推理证明。答案正确可能来自巧合，测试通过也可能过拟合。形式证明由核验器检查每个推导步骤，具有更强的逻辑证书，但自动形式化及题意对应又是额外问题。AlphaProof 译文把这一点具体化；不能把形式核验结论扩展到任意自然语言研究。

R1 正式论文比较无冷启动的 R1-Zero 与多阶段 R1，并报告强化学习中轨迹长度、自我检查等现象。它们是可观察文本和性能现象，不是意识或全部内部推理过程的直接记录。多阶段流程同时改变数据、训练目标与模型状态；一个最后榜单分数不能识别每项的独立因果贡献，应读对应消融、评价样本与采样预算。

\originalsection{进步来自新能力还是选择概率}
把基础模型在一道题上生成正确解答的概率记为 $p$。若独立采样 $k$ 次、存在完美选择器，则 pass@$k$ 为 $1-(1-p)^k$。训练把 $p$ 从 $0.05$ 提到 $0.30$ 时，单次成功大幅提高；足够大的基础模型采样预算也可能暴露已有的解题路径。这两个事实不冲突。

关于 RLVR 能力边界的研究\cite{rlvrboundary}使用 pass@$k$、覆盖与生成似然比较基础和训练后分布。其意义是提醒我们同时测高概率能力和可达的正确轨迹；有限 $k$ 没有采到某轨迹，不能证明其概率严格为零。反过来，基础模型偶尔能生成正确答案，也不能说训练没有实际收益。学习者应把“可达范围”“固定预算可靠性”和“跨分布迁移”作为三个不同目标。

本项目的有限动作奖励实验实际更新指定策略，并记录了错误验证器下 best-of-$N$ 随预算增大反而降低真正确率的案例。它检验奖励和选择偏差的数学机制，没有训练 671B R1，更没有用微型实验声称复现其长链推理。下一次可行研究是固定小模型和题库，比较更新前后 pass@1、固定 $k$ 的覆盖、输出长度和分布外题目；测试题不可参与挑模型或奖励规则。

\originalsection{练习与下一步阅读}
阅读顺序是教材的后训练推导、DPO 全译的闭式策略变换、R1 Methods 与实验、Snell 全译的预算分配，再到 CaT 的无参考监督。每篇都问：奖励从哪里来？哪个对象被优化？选择器能访问哪些信息？预算与比较对象是否匹配？这些问题比仅记住方法缩写更能帮助独立设计研究。

\begin{exercise}
若 $p=0.1$，完美验证器下十次独立采样成功概率是多少？若所有十次输出完全相同呢？
\end{exercise}
\begin{solution}
独立时为 $1-0.9^{10}\approx0.6513$；完全相同时仍为 $0.1$。多样性和相关性是预算收益的必要测量，不能只报采样次数。
\end{solution}
\begin{exercise}
训练后的正确回答更长，能否说明长回答导致正确？如何得到更有用的证据？
\end{exercise}
\begin{solution}
不能；难题、模型状态、训练筛选可能同时影响长度与正确性。应在可控任务上固定模型和输入，改变可用预算或训练条件，报告成对结果及截断如何影响内容；还要区别真正额外计算与无效重复文本。
\end{solution}
